\SourceLeaf{387}%
\[
\theta=-\frac{A'A''+B'B''+C'C''}{n},
\]
also [ist]
\[
\Sigma\alpha(cB'-bC')=0,
\]
folglich auch dessen Differential $=0$. Ferner
\[
\begin{aligned}
\Sigma a''(cB'-bC')&=\theta(Aa''+Bb''+Cc'')\\
&=-\theta(aA''+bB''+cC'')\\
&=+\theta\theta.
\end{aligned}
\]
\par\indent Wir haben folglich:
\[
\begin{aligned}
mm\left(\frac{\partial\partial\log m}{\partial t^2}+\frac{\partial\partial\log m}{\partial u^2}\right)
&=-(aA'+bB'+cC')(\alpha A''+\beta B''+\gamma C'')\\
&\quad+\frac{(A'A''+B'B''+C'C'')^2}{nn}.
\end{aligned}
\]
\par\indent Setzt man also
\[
A(B'C''-C'B'')+B(C'A''-A'C'')+C(A'B''-B'A'')=\Delta,
\]
so wird
\[
\begin{aligned}
&mm\left(\frac{\partial\partial\log m}{\partial t^2}+\frac{\partial\partial\log m}{\partial u^2}\right)\\
&\quad=-\frac{\begin{gathered}(A'A'+B'B'+C'C'+\Delta)(A''A''+B''B''+C''C''+\Delta)\\[-2pt]{}-(A'A''+B'B''+C'C'')^2\end{gathered}}{nn},
\end{aligned}
\]
[oder]
\begin{equation}\tag*{[24]}
\frac{\partial\partial\log m}{\partial t^2}+\frac{\partial\partial\log m}{\partial u^2}=-\Delta.
\end{equation}
Wir bemerken noch, dass
\begin{equation}\tag*{$(\FourRingCross)$}
\left\{\begin{aligned}
B'C''-C'B''&=A\Delta\\
C'A''-A'C''&=B\Delta\\
A'B''-B'A''&=C\Delta.
\end{aligned}\right.
\end{equation}
\par\indent Auch haben wir uns überzeugt, dass in den Bedingungsgleichungen
\[
\frac{\partial a'}{\partial u}=\frac{\partial\alpha'}{\partial t},\quad\text{etc.}
\]
nichts weiter liegt als dieses Resultat [nemlich die Gleichung 24].

[Die Gleichung 24 hat folgende geometrische Bedeutung. Wird das Flächenelement
\[
mm\,dt\,du
\]
